% =====================================================================
\begin{thebibliography}{99}
\addcontentsline{toc}{section}{References}

\bibitem{ABD07} A.~Atserias, A.~Bulatov and V.~Dalmau. On the power of $k$-consistency. In \emph{Automata, Languages and Programming (ICALP 2007)}, LNCS 4596, Springer (2007), 279--290.

\bibitem{BartoKozik14} L.~Barto and M.~Kozik. Constraint satisfaction problems solvable by local consistency methods. \emph{J.\ ACM} 61 (2014), no.~1, Article~3. \href{https://doi.org/10.1145/2556646}{doi:10.1145/2556646}.

\bibitem{Berger66} R.~Berger. The undecidability of the domino problem. \emph{Mem.\ Amer.\ Math.\ Soc.} 66 (1966). \href{https://doi.org/10.1090/memo/0066}{doi:10.1090/memo/0066}.

\bibitem{BermanSimon88} P.~Berman and J.~Simon. Investigations of fault-tolerant networks of computers. In \emph{Proceedings of STOC 1988}, ACM (1988), 66--77. \href{https://doi.org/10.1145/62212.62219}{doi:10.1145/62212.62219}.

\bibitem{BMM13} A.~Bu\v{s}i\'c, J.~Mairesse and I.~Marcovici. Probabilistic cellular automata, invariant measures, and perfect sampling. \emph{Adv.\ Appl.\ Probab.} 45 (2013), 960--980.

\bibitem{Choquet} G.~Choquet. Existence et unicit\'e des repr\'esentations int\'egrales au moyen des points extr\'emaux dans les c\^ones convexes. \emph{S\'eminaire Bourbaki}, expos\'e 139 (1956--1958), 33--47.

\bibitem{CPY89} K.~\v{C}ulik~II, J.~Pachl and S.~Yu. On the limit sets of cellular automata. \emph{SIAM J.\ Comput.} 18 (1989), 831--842. \href{https://doi.org/10.1137/0218057}{doi:10.1137/0218057}.

\bibitem{DKV02} V.~Dalmau, Ph.~G.~Kolaitis and M.~Y.~Vardi. Constraint satisfaction, bounded treewidth, and finite-variable logics. In \emph{Principles and Practice of Constraint Programming (CP 2002)}, LNCS 2470, Springer (2002), 310--326.

\bibitem{DM26} S.~Douglas and N.~Mghirbi. Exponential ergodicity and the space--time Gibbs states of one-dimensional probabilistic cellular automata. Preprint, version 1.2.0, Zenodo (2026). \href{https://doi.org/10.5281/zenodo.24122307}{doi:10.5281/zenodo.24122307}. Numbers of results refer to this version.

\bibitem{DMforgetting14} S.~Douglas and N.~Mghirbi. Forgetting under noise is undecidable. Preprint, version 1.4.0, Zenodo (2026). \href{https://doi.org/10.5281/zenodo.24126025}{doi:10.5281/zenodo.24126025}. Machine-family Lean core.

\bibitem{DMnil} S.~Douglas and N.~Mghirbi. A nilpotent cellular automaton that remembers under noise. Preprint, version 1.2.0, Zenodo (2026). \href{https://doi.org/10.5281/zenodo.24120279}{doi:10.5281/zenodo.24120279}.

\bibitem{DMdiagram} S.~Douglas and N.~Mghirbi. One noisy cellular automaton with an uncomputable noise diagram. Preprint, version 1.1.0, Zenodo (2026). \href{https://doi.org/10.5281/zenodo.24152233}{doi:10.5281/zenodo.24152233}. All numbered locators refer to this version.
\bibitem{DMdiagramLean} S.~Douglas and N.~Mghirbi. One noisy cellular automaton with an uncomputable noise diagram. Preprint, version 1.2.0, with its Lean formalization, Zenodo (2026). \href{https://doi.org/10.5281/zenodo.24174445}{doi:10.5281/zenodo.24174445}.

\bibitem{DMhard} S.~Douglas and N.~Mghirbi. An uncomputable phase diagram at positive temperature for one classical lattice model. Preprint, version 1.2.0, Zenodo (2026). \href{https://doi.org/10.5281/zenodo.24144192}{doi:10.5281/\allowbreak zenodo.24144192}.

\bibitem{Dob68} R.~L.~Dobrushin. The description of a random field by means of conditional probabilities and conditions of its regularity. \emph{Theory Probab.\ Appl.} 13 (1968), 197--224.

\bibitem{DS85} R.~L.~Dobrushin and S.~B.~Shlosman. Constructive criterion for the uniqueness of Gibbs field. In \emph{Statistical Physics and Dynamical Systems}, Progress in Physics 10, Birkh\"auser (1985), 347--370.

\bibitem{FederVardi98} T.~Feder and M.~Y.~Vardi. The computational structure of monotone monadic SNP and constraint satisfaction: a study through Datalog and group theory. \emph{SIAM J.\ Comput.} 28 (1998), 57--104. \href{https://doi.org/10.1137/S0097539794266766}{doi:10.1137/S0097539794266766}.

\bibitem{Gacs01} P.~G\'acs. Reliable cellular automata with self-organization. \emph{J.\ Stat.\ Phys.} 103 (2001), 45--267; revised version arXiv:math/0003117v2 (2024).

\bibitem{Gacs21} P.~G\'acs. A new version of Toom's proof. Preprint, arXiv:2105.05968 (2021).

\bibitem{GacsReif88} P.~G\'acs and J.~Reif. A simple three-dimensional real-time reliable cellular array. \emph{J.\ Comput.\ System Sci.} 36 (1988), 125--147. \href{https://doi.org/10.1016/0022-0000(88)90024-4}{doi:10.1016/0022-0000(88)90024-4}.

\bibitem{GKLM89} S.~Goldstein, R.~Kuik, J.~L.~Lebowitz and C.~Maes. From PCA's to equilibrium systems and back. \emph{Comm.\ Math.\ Phys.} 125 (1989), 71--79.

\bibitem{Halawi22} A.~Halawi. \emph{Asymptotic behaviour of probabilistic cellular automata}. MS thesis, American University of Beirut (2022). \url{https://scholarworks.aub.edu.lb/handle/10938/23549}. Theorem~4.2.3 and Proposition~4.2.4.

\bibitem{Helly23} E.~Helly. \"Uber Mengen konvexer K\"orper mit gemeinschaftlichen Punkten. \emph{Jahresber.\ Deutsch.\ Math.-Verein.} 32 (1923), 175--176.

\bibitem{Hoeffding63} W.~Hoeffding. Probability inequalities for sums of bounded random variables. \emph{J.\ Amer.\ Statist.\ Assoc.} 58 (1963), 13--30. \href{https://doi.org/10.1080/01621459.1963.10500830}{doi:10.1080/01621459.1963.10500830}.

\bibitem{JCG97} P.~Jeavons, D.~Cohen and M.~Gyssens. Closure properties of constraints. \emph{J.\ ACM} 44 (1997), 527--548.

\bibitem{Kari92} J.~Kari. The nilpotency problem of one-dimensional cellular automata. \emph{SIAM J.\ Comput.} 21 (1992), 571--586. \href{https://doi.org/10.1137/0221036}{doi:10.1137/0221036}.

\bibitem{Kechris} A.~S.~Kechris. \emph{Classical Descriptive Set Theory}. Graduate Texts in Mathematics 156, Springer (1995).

\bibitem{KolaitisVardi00} Ph.~G.~Kolaitis and M.~Y.~Vardi. A game-theoretic approach to constraint satisfaction. In \emph{Proceedings of AAAI 2000}, AAAI Press (2000), 175--181.

\bibitem{Kurdyumov78} G.~L.~Kurdyumov. An algorithm-theoretic method for the study of homogeneous random networks. In \emph{Locally Interacting Systems and Their Application in Biology}, Lecture Notes in Math.\ 653, Springer (1978), 37--55.

\bibitem{Lean4} L.~de~Moura and S.~Ullrich. The Lean~4 theorem prover and programming language. In \emph{Automated Deduction -- CADE 28}, LNCS 12699, Springer (2021), 625--635.

\bibitem{LeanWang} \emph{lean-wang}: a Lean formalization of results on Wang tilings, GitHub repository \url{https://github.com/edemaine/lean-wang}, directory \texttt{LeanWang/Robinson/Closed104}.

\bibitem{LMS90} J.~L.~Lebowitz, C.~Maes and E.~R.~Speer. Statistical mechanics of probabilistic cellular automata. \emph{J.\ Stat.\ Phys.} 59 (1990), 117--170.

\bibitem{Mathlib20} The mathlib Community. The Lean mathematical library. In \emph{Proceedings of the 9th ACM SIGPLAN International Conference on Certified Programs and Proofs (CPP 2020)}, ACM (2020), 367--381.

\bibitem{MST19} I.~Marcovici, M.~Sablik and S.~Taati. Ergodicity of some classes of cellular automata subject to noise. \emph{Electron.\ J.\ Probab.} 24 (2019), paper no.~41, 1--44. \href{https://doi.org/10.1214/19-EJP297}{doi:10.1214/19-EJP297}.

\bibitem{MT26} I.~Marcovici and S.~Taati. Positive-rate PCA and IPS with stationary Bernoulli measures are rapidly forgetful. Preprint, arXiv:2605.16904 (2026).

\bibitem{MarsanSablikTorma26} H.~Marsan, M.~Sablik and I.~T\"orm\"a. A perturbed cellular automaton with two phase transitions for the ergodicity. \emph{J.\ Stat.\ Phys.} 193 (2026), 6. \href{https://doi.org/10.1007/s10955-025-03564-0}{doi:10.1007/s10955-025-03564-0}.

\bibitem{MiekiszTaati26} J.~Mi\k{e}kisz and S.~Taati. Quasicrystalline Gibbs states in 4-dimensional lattice-gas models with finite-range interactions. \emph{J.\ Stat.\ Phys.} 193 (2026), 166. \href{https://doi.org/10.1007/s10955-026-03727-7}{doi:10.1007/s10955-026-03727-7}.

\bibitem{Mosse92} B.~Moss\'e. Puissances de mots et reconnaissabilit\'e des points fixes d'une substitution. \emph{Theoret.\ Comput.\ Sci.} 99 (1992), 327--334.

\bibitem{Mozes89} S.~Mozes. Tilings, substitution systems and dynamical systems generated by them. \emph{J.\ Analyse Math.} 53 (1989), 139--186.

\bibitem{Ollinger08} N.~Ollinger. Two-by-two substitution systems and the undecidability of the domino problem. In \emph{Logic and Theory of Algorithms (CiE 2008)}, LNCS 5028, Springer (2008), 476--485.

\bibitem{PW96} J.~G.~Propp and D.~B.~Wilson. Exact sampling with coupled Markov chains and applications to statistical mechanics. \emph{Random Structures Algorithms} 9 (1996), 223--252.

\bibitem{Radon21} J.~Radon. Mengen konvexer K\"orper, die einen gemeinsamen Punkt enthalten. \emph{Math.\ Ann.} 83 (1921), 113--115. \href{https://doi.org/10.1007/BF01464231}{doi:10.1007/BF01464231}.

\bibitem{Robinson71} R.~M.~Robinson. Undecidability and nonperiodicity for tilings of the plane. \emph{Invent.\ Math.} 12 (1971), 177--209.

\bibitem{Sla81} J.~Slawny. Ergodic properties of equilibrium states. \emph{Comm.\ Math.\ Phys.} 80 (1981), 477--483.

\bibitem{Solomyak98} B.~Solomyak. Nonperiodicity implies unique composition for self-similar translationally finite tilings. \emph{Discrete Comput.\ Geom.} 20 (1998), 265--279.

\bibitem{SST26} J.~M.~Swart, R.~Szab\'o and C.~Toninelli. Peierls bounds from Toom contours. \emph{J.\ Theor.\ Probab.} (2026). \href{https://doi.org/10.1007/s10959-026-01508-2}{doi:10.1007/s10959-026-01508-2}; arXiv:2202.10999.

\bibitem{Toom80} A.~L.~Toom. Stable and attractive trajectories in multicomponent systems. In \emph{Multicomponent Random Systems}, R.~L.~Dobrushin and Ya.~G.~Sinai (eds.), Adv.\ Probab.\ Related Topics 6, Dekker, New York (1980), 549--575.

\bibitem{Toom00} A.~Toom. Algorithmic unsolvability of the ergodicity problem for binary cellular automata. \emph{Markov Process.\ Related Fields} 6 (2000), 569--577.

\bibitem{vdBS99} J.~van den Berg and J.~E.~Steif. On the existence and nonexistence of finitary codings for a class of random fields. \emph{Ann.\ Probab.} 27 (1999), 1501--1522. Theorem numbers as in the preprint version.

\bibitem{Wang61} H.~Wang. Proving theorems by pattern recognition~II. \emph{Bell System Tech.\ J.} 40 (1961), 1--41.
\bibitem{Kari05} J.~Kari. Theory of cellular automata: A survey. \emph{Theoret.\ Comput.\ Sci.} 334 (2005), 3--33. \href{https://doi.org/10.1016/j.tcs.2004.11.021}{doi:10.1016/j.tcs.2004.11.021}.

\end{thebibliography}




